\documentclass[12pt,a4paper]{article}

% ── Pacotes ────────────────────────────────────────────────────────────────────
\usepackage[utf8]{inputenc}
\usepackage[T1]{fontenc}
\usepackage[english]{babel}
\usepackage{amsmath,amssymb,amsthm}
\usepackage{mathrsfs}
\usepackage{geometry}
\usepackage{xcolor}
\usepackage{hyperref}
\usepackage{booktabs}
\usepackage{array}
\usepackage{enumitem}

\geometry{
	a4paper,
	top=2.5cm, bottom=2.5cm,
	left=3cm, right=3cm
}

\hypersetup{
	colorlinks=true,
	linkcolor=blue,
	citecolor=blue,
	urlcolor=blue,
	pdftitle={Reformulacao Espectral do Motor de Heranca Estrutural},
	pdfauthor={Tiago Bandeira}
}

% ── Ambientes de Teorema ───────────────────────────────────────────────────────
\newtheorem{theorem}{Teorema}[section]
\newtheorem{lemma}[theorem]{Lema}
\newtheorem{proposition}[theorem]{Proposição}
\newtheorem{corollary}[theorem]{Corolário}

\theoremstyle{definition}
\newtheorem{definition}[theorem]{Definição}
\newtheorem{example}[theorem]{Exemplo}
\newtheorem{remark}[theorem]{Observação}

% ── Atalhos matemáticos ────────────────────────────────────────────────────────
\newcommand{\N}{\mathbb{N}}
\newcommand{\Z}{\mathbb{Z}}
\newcommand{\R}{\mathbb{R}}
\newcommand{\C}{\mathbb{C}}
\newcommand{\T}{\mathbb{T}}
\newcommand{\PP}{\mathcal{P}}          % conjunto dos primos
\newcommand{\F}{\mathcal{F}}           % fita-dobra
\newcommand{\G}{\mathcal{G}}           % grade
\newcommand{\LL}{\mathcal{L}}
\newcommand{\HRm}{\mathrm{HR}^{-}}
\newcommand{\HRp}{\mathrm{HR}^{+}}
\newcommand{\Lam}{\Lambda}             % von Mangoldt
\newcommand{\e}[1]{e^{2\pi i (#1)}}
\newcommand{\norm}[1]{\left\| #1 \right\|}
\newcommand{\abs}[1]{\left| #1 \right|}
\newcommand{\ip}[2]{\langle #1,\, #2 \rangle}
\DeclareMathOperator{\Spec}{Spec}

% ── Título ─────────────────────────────────────────────────────────────────────
\title{%
	\textbf{Reformulação Espectral do Motor de Herança Estrutural}\\[0.4em]
	\large Uma Abordagem por Sistemas Dinâmicos à Conjectura de Goldbach
}
\author{Tiago Bandeira\\[0.3em]
	\small\texttt{tiagobandeira.goldbach2025@gmail.com}
}
\date{Maio de 2026 \\ \vspace{0.3em}
	\small Preprint — Artigo 6 da série sobre a Conjectura de Goldbach}

% ═══════════════════════════════════════════════════════════════════════════════
\begin{document}
	
	\maketitle
	\thispagestyle{empty}
	
	% ── Nota do Autor ─────────────────────────────────────────────────────────────
	\begin{center}
		\fbox{\parbox{0.88\textwidth}{
				\textbf{Nota do Autor.} Este artigo não reivindica uma demonstração da
				Conjectura de Goldbach. Seu objetivo é reformular a questão central —
				a Hipótese Restritiva Fraca $\HRm$ — na linguagem de sistemas dinâmicos
				e análise espectral, e provar incondicionalmente que toda a interferência
				espectral das frequências não nulas se cancela (via teorema de
				Green-Tao-Ziegler \cite{greentaoziegler}), reduzindo o problema a uma
				única componente de frequência zero. Os resultados incondicionais são
				explicitamente identificados; os condicionais são declarados como tais.
				Comentários e críticas matemáticas são bem-vindos.
		}}
	\end{center}
	\vspace{0.5em}
	
	% ── Abstract ──────────────────────────────────────────────────────────────────
	\begin{abstract}
		Reformulamos o problema central do Motor de Herança Estrutural
		\cite{motor} na linguagem de sistemas dinâmicos e análise espectral.
		A Hipótese Restritiva Fraca ($\HRm$) — existência de um par de primos numa janela $W_i$ da
		enumeração espacial $\sigma$ — é traduzida, de forma
		incondicional, para a condição $\widehat{F}(0) > 0$: positividade da
		componente de frequência zero da função indicadora de colisão
		$F(i) = \mathbf{1}_{\PP}(e_1(i))\cdot\mathbf{1}_{\PP}(e_3(i))$ no
		sistema dinâmico de medida finito $(X_N, \mu_N, T)$.
		
		Demonstramos incondicionalmente que o operador de Koopman $U_T$ possui
		espectro discreto puro (raízes $K$-ésimas da unidade, $K = N-1$), que
		$T$ é ergódico com entropia de Kolmogorov-Sinai nula, e que $\HRm
		\Leftrightarrow \widehat{F}(0) > 0$. Adicionalmente, provamos de forma
		incondicional — via o teorema de Green-Tao-Ziegler
		\cite{greentaoziegler} sobre ortogonalidade da função de von Mangoldt
		a nilsequências — que $\widehat{w}(j) \to 0$ para todo $j \neq 0$.
		O problema se reduz assim, incondicionalmente, a uma única componente
		espectral: a positividade de $\widehat{w}(0)$, que é equivalente a
		$\HRm$ e permanece em aberto.
		
		Discutimos as razões precisas pelas quais os métodos clássicos — método
		do círculo de Hardy-Littlewood-Vinogradov e crivo linear de
		Rosser-Iwaniec — colapsam nesta configuração, e delineamos como a
		formulação espectral reposiciona a barreira analítica remanescente em
		termos acessíveis às técnicas modernas.
	\end{abstract}
	
	\tableofcontents
	\newpage
	
	% ═══════════════════════════════════════════════════════════════════════════════
	\section{Introdução}
	\label{sec:intro}
	
	\subsection{A série e o contexto}
	
	Este artigo integra a série \cite{goldbach1,grade,ancora,crivo_b,motor}
	dedicada a construir um mecanismo dinâmico de cobertura dos números pares
	pela Conjectura de Goldbach. O Motor de Herança Estrutural \cite{motor}
	reduz o problema global
	\begin{quote}
		\emph{``Todo par $2M > 2$ é soma de dois primos''}
	\end{quote}
	à questão local, repetida para cada $M$:
	\begin{quote}
		\emph{``Existe algum passo $t \in \{1,\ldots,N-1\}$ da órbita de
			$\sigma$ tal que $e_1(t)$ e $e_3(t)$ são ambos primos?''}
	\end{quote}
	Esta é precisamente $\HRm$. A camada algébrica do motor — bijeção $\Phi$,
	cobertura pela órbita de $\sigma$, mecanismo de promoção — foi demonstrada
	incondicionalmente em \cite{motor}. A única lacuna remanescente é a
	primalidade simultânea dos extremos de $W_{i^*}$ em algum passo da órbita.
	
	O presente artigo não resolve essa lacuna. Seu propósito é reescrevê-la
	na linguagem de sistemas dinâmicos e análise harmônica, tornando explícita
	a natureza da obstrução e alinhando o problema com a escola moderna de
	teoria analítica dos números.
	
	\subsection{Por que os métodos clássicos falham}
	
	As três principais tentativas de provar $\HRm$ diretamente pela teoria
	analítica dos números clássica colapsam por razões estruturais precisas,
	detalhadas na Seção~\ref{sec:crivo}:
	
	\begin{enumerate}[label=(\roman*)]
		\item \textbf{Método do círculo (Hardy-Littlewood-Vinogradov).} A
		separação de escalas — comprimento de progressão $L \sim (\log M^-)^2$
		versus alvo $2M^- \gg L$ — destrói o termo principal da integral de
		Fourier quando $M \to \infty$.
		
		\item \textbf{Crivo linear de Rosser-Iwaniec.} O erro acumulado
		$\mathcal{R}$ cresce como $(\log M^-)^{3-\varepsilon}$, dominando o
		termo principal de ordem $(\log M^-)^2 / (\log\log M^-)^2$.
		
		\item \textbf{Indução algébrica pura.} A primalidade é multiplicativa;
		a fita-dobra opera sob leis aditivas. Transportar primalidade passo a
		passo cai em circularidade indutiva.
	\end{enumerate}
	
	O diagnóstico comum é claro: \emph{ferramentas estáticas de ponto fixo não
		resolvem um sistema dinâmico em evolução}. A reformulação proposta analisa
	o fluxo orbital global em vez de instâncias isoladas.
	
	\subsection{Mudança de perspectiva: da contagem ao fluxo}
	
	A filosofia central é a seguinte transição de linguagem:
	
	\begin{center}
		\renewcommand{\arraystretch}{1.4}
		\begin{tabular}{lll}
			\toprule
			\textbf{Dimensão} & \textbf{Abordagem Clássica} & \textbf{Abordagem Espectral} \\
			\midrule
			Análise & Pontual (um $2M$ de cada vez) & Global (fluxo sobre $X_N$) \\
			Matemática & Aritmética de contagem & Geometria de medida \\
			O motor & Organizador estático & Operador espectral (Koopman) \\
			A barreira & Paridade e acúmulo de $\mathcal{R}$ & Freq.\ zero de $\widehat{w}$ (GTZ isola) \\
			\bottomrule
		\end{tabular}
	\end{center}
	
	Esta transição sintoniza o Motor de Herança com a escola moderna de
	Teoria Analítica dos Números (Green-Tao-Gowers \cite{greentao},
	Green-Tao-Ziegler \cite{greentaoziegler}, Sarnak \cite{sarnak}),
	na qual a pseudo-aleatoriedade dos primos é explorada como ferramenta,
	não como obstáculo.
	
	\begin{remark}[O que a reformulação oferece e o que não oferece]
		\label{rem:honestidade}
		A equivalência $\HRm \Leftrightarrow \widehat{F}(0) > 0$
		(Teorema~\ref{thm:reformulacao}) é uma tautologia precisa, não uma
		prova. Ela não reduz a dificuldade de Goldbach; ela a \emph{localiza}
		numa obstrução espectral única e bem definida. O valor está em que essa
		obstrução é exatamente o tipo de objeto para o qual existem ferramentas
		modernas — nilsequências, normas de Gowers, princípio de Sarnak — que
		ainda não estavam disponíveis nas formulações clássicas.
	\end{remark}
	
	\subsection{Organização do artigo}
	
	A Seção~\ref{sec:fundamentos} recapitula os fundamentos algébricos do
	motor. A Seção~\ref{sec:dinamica} constrói o sistema dinâmico
	$(X_N, \mu_N, T)$ e prova o Teorema Espectral (incondicional).
	A Seção~\ref{sec:colisao} define a função de colisão $F$ e estabelece
	$\HRm \Leftrightarrow \widehat{F}(0) > 0$ (incondicional).
	A Seção~\ref{sec:obstrucao} descreve a obstrução espectral e sua conexão
	com Sarnak e Gowers. A Seção~\ref{sec:crivo} detalha o colapso dos
	métodos clássicos. A Seção~\ref{sec:conclusao} sintetiza o alcance e
	aponta direções futuras.
	
	% ═══════════════════════════════════════════════════════════════════════════════
	\section{Fundamentos Algébricos do Motor}
	\label{sec:fundamentos}
	
	Nesta seção resumimos os objetos e resultados incondicionais de
	\cite{motor} necessários para a formulação espectral. Todos os resultados
	desta seção são independentes de primalidade.
	
	\subsection{A fita-dobra e a grade}
	
	\begin{definition}[Fita-dobra $\F_N^\varepsilon$]
		\label{def:fita}
		Fixados $N \geq 1$ e $\varepsilon \in \{0,1\}$, a fita-dobra
		$\F_N^\varepsilon$ é a matriz $2 \times N$ com entradas
		\[
		a_k = 2k - 1,\qquad b_k = 4N - 2k + 1 - 2\varepsilon,\qquad k = 1,\ldots,N.
		\]
		Para todo $k$ e $\varepsilon$, $a_k + b_k = 4N - 2\varepsilon =: 2M(\varepsilon)$.
	\end{definition}
	
	\begin{definition}[Grade $G_{3,C}$]
		Para $C \geq 2$ par, a grade $G_{3,C}$ é a matriz $3 \times C$ cujas
		células contêm os primeiros $3C$ ímpares positivos, organizados por linhas:
		\[
		f(r,c) = 2\bigl[(r-1)C + c\bigr] - 1,\quad r \in \{1,2,3\},\; c \in \{1,\ldots,C\}.
		\]
	\end{definition}
	
	\begin{definition}[Acoplamento $\Phi$]
		Sob a condição $3C = 2N$, o acoplamento $\Phi: \{1,\ldots,N\} \to
		G_{3,C}\times G_{3,C}$ associa a cada coluna $k$ da fita um par de células
		de $G_{3,C}$, preservando a soma $\Phi(k)_1 + \Phi(k)_2 = 4N = 2M^+$
		para todo $k$ \cite[Prop.~3.5]{motor}.
	\end{definition}
	
	\subsection{A permutação deslizante e sua órbita}
	
	\begin{definition}[Permutação deslizante $\sigma$]
		Seja $S$ o conjunto dos $3C-2$ elementos interiores do percurso zigzag
		de $G_{3,C}$. A permutação deslizante $\sigma: S \to S$ avança cada
		elemento uma posição no percurso zigzag: $\sigma(\sigma_G(k)) =
		\sigma_G(k-1)$.
	\end{definition}
	
	A janela de índice $i^*$ (com $i^* = \lfloor C/2 \rfloor$) percorre,
	sob a enumeração espacial $\sigma$, exatamente os $N-1$ pares antipodais
	$(e_1(i), e_3(i))$ com $e_1(i) + e_3(i) = 2M^+ = 4N$ — resultado
	incondicional, corrigido de \cite[Prop.~4.3, Prop.~4.6]{motor}.
	
	\begin{remark}[Correcção da soma constante]
		\label{rem:correcao-soma}
		Versões anteriores deste artigo e de \cite{motor} afirmavam que a
		soma constante dos extremos da janela era $2M^- = 4N-2$.  Esta
		afirmação estava incorrecta.  A Proposição~3.5 de \cite{motor} já
		prova que $\Phi(k)_1 + \Phi(k)_2 = 4N = 2M^+$ para todo $k$.
		A correcção é que os extremos $(e_1(i), e_3(i))$ de $W_i$
		satisfazem $e_1(i)+e_3(i) = 2M^+$, não $2M^-$.  A grade aponta
		geometricamente para o alvo futuro; nenhum incremento adicional
		é necessário.
	\end{remark}
	
	\begin{remark}[Equações de indexação espacial]
		\label{rem:movimento}
		Por construção do percurso zigzag e do acoplamento $\Phi$, os
		extremos de $W_i$ variam linearmente com o índice espacial $i$:
		\[
		e_1(i) = e_1(1) + 2(i-1),\qquad e_3(i) = e_3(1) - 2(i-1),
		\]
		de modo que $e_1(i) + e_3(i) = e_1(1) + e_3(1) = 2M^+$ é
		constante para todo $i \in \{1,\ldots,N-1\}$.
		
		\textbf{Correcção:} a soma constante é $2M^+ = 4N$,
		não $2M^- = 4N-2$.  O índice $i$ é \textbf{espacial} — identifica
		qual janela (par antipodal) estamos a observar na grade fixa — e
		não uma iteração temporal.  A grade é um ``oráculo'' geométrico do
		alvo futuro $2M^+$.
	\end{remark}
	
	\subsection{As hipóteses restritivas}
	
	\begin{definition}[Hipóteses Restritivas]
		\label{def:HR}
		Fixados $C$ e $i^* = \lfloor C/2 \rfloor$, sejam $(e_1(i), e_3(i))$ os
		extremos da janela $W_i$, com $e_1(i)+e_3(i)=2M^+$ por construção
		geométrica.
		
		\begin{itemize}
			\item \textbf{$\HRm$ (Fraca):} $\exists\, i \in \{1,\ldots,N-1\}$
			tal que $e_1(i) \in \PP$ e $e_3(i) \in \PP$.
			\item \textbf{$\HRp$ (Forte):} $\exists\, i \in \{1,\ldots,N-1\}$
			tal que a tripla $(e_1(i), p_2(i), e_3(i))$ de $W_i$ é
			inteiramente formada por primos.
		\end{itemize}
	\end{definition}
	
	Sob $\HRm$, o Teorema do Motor \cite[Thm.~7.3]{motor} garante que a
	cadeia $2M_0 \to 2M_0 + 2 \to \cdots$ cobre todos os pares $\geq 2M_0$
	como somas de dois primos. A lacuna central — a única hipótese não provada
	— é precisamente $\HRm$.
	
	% ═══════════════════════════════════════════════════════════════════════════════
	\section{O Sistema Dinâmico de Medida Finito}
	\label{sec:dinamica}
	
	\subsection{Construção do sistema}
	
	\begin{definition}[Espaço de estados $X_N$]
		\label{def:estados}
		O espaço de estados do motor é o conjunto finito de índices
		\textbf{espaciais}
		\[
		X_N := \{i \in \Z : 1 \leq i \leq N-1\}.
		\]
		Cada estado $i \in X_N$ corresponde biunivocamente ao par antipodal
		$(e_1(i), e_3(i))$ da grade, com soma $e_1(i)+e_3(i)=2M^+$
		garantida geometricamente.  O índice $i$ identifica
		\emph{qual janela} $W_i$ estamos a observar, não um instante de
		tempo.
	\end{definition}
	
	\begin{definition}[Medida de probabilidade uniforme $\mu_N$]
		Definimos a medida de probabilidade sobre a $\sigma$-álgebra discreta
		de $X_N$ por
		\[
		\mu_N(\{i\}) = \frac{1}{N-1},\qquad \forall\, i \in X_N.
		\]
	\end{definition}
	
	\begin{definition}[Transformação dinâmica $T$]
		A transformação $T: X_N \to X_N$ é a translação cíclica unitária
		\[
		T(i) = i + 1 \pmod{N-1}.
		\]
		$T$ percorre ciclicamente todos os índices espaciais das janelas,
		equivalendo a observar a janela seguinte na grade.
	\end{definition}
	
	\begin{lemma}[Preservação de medida e ergodicidade]
		\label{lem:ergodico}
		A transformação $T$ é um automorfismo de medida sobre $(X_N, \mu_N)$.
		Além disso, $(X_N, \mu_N, T)$ é ergódico e possui entropia métrica de
		Kolmogorov-Sinai nula:
		\[
		\mu_N(T^{-1}(B)) = \mu_N(B)\;\;\forall B \subset X_N,\qquad h(T) = 0.
		\]
	\end{lemma}
	
	\begin{proof}
		Como $T$ é uma bijeção de $X_N$ sobre si mesmo e $\mu_N$ é a medida
		de contagem normalizada, a preservação de medida é imediata. A
		ergodicidade decorre do fato de que $T$ forma um único ciclo de
		comprimento máximo $K = N-1$ sobre $X_N$: o único subconjunto $B$ com
		$T^{-1}(B) = B$ é $B = \emptyset$ ou $B = X_N$. A entropia nula é
		consequência do caráter periódico e determinístico de $T$.
	\end{proof}
	
	\subsection{O operador de Koopman e sua estrutura espectral}
	
	\begin{definition}[Espaço de Hilbert e produto interno]
		Seja $L^2(X_N, \mu_N)$ o espaço de funções de valor complexo sobre
		$X_N$, com produto interno
		\[
		\ip{\psi}{\phi} = \frac{1}{N-1}\sum_{i=1}^{N-1} \psi(i)\,\overline{\phi(i)}.
		\]
	\end{definition}
	
	\begin{definition}[Operador de Koopman $U_T$]
		\label{def:koopman}
		O operador de Koopman $U_T: L^2(X_N) \to L^2(X_N)$ é definido por
		\[
		(U_T\psi)(i) = \psi(T(i)).
		\]
		Como $T$ preserva $\mu_N$, $U_T$ é unitário: $U_T^* = U_T^{-1}$.
		Geometricamente, $U_T$ translada a janela de observação sobre a
		grade fixa: em vez de mover os números, move o ``enquadramento''.
	\end{definition}
	
	\begin{proposition}[Representação matricial circulante]
		\label{prop:circulante}
		Na base canônica de $L^2(X_N)$, o operador $U_T$ é representado pela
		matriz de permutação circulante $M \in \R^{K\times K}$ (com $K = N-1$):
		\[
		M = \begin{pmatrix}
			0 & 1 & 0 & \cdots & 0 \\
			0 & 0 & 1 & \cdots & 0 \\
			\vdots & \vdots & \vdots & \ddots & \vdots \\
			0 & 0 & 0 & \cdots & 1 \\
			1 & 0 & 0 & \cdots & 0
		\end{pmatrix}.
		\]
	\end{proposition}
	
	\begin{theorem}[Espectro discreto do motor — incondicional]
		\label{thm:espectro}
		O espectro de $U_T$ é discreto puro. Os autovalores e as autofunções
		associadas são:
		\[
		\lambda_j = \exp\!\left(\frac{2\pi i\, j}{K}\right),\qquad
		\phi_j(i) = \exp\!\left(\frac{2\pi i\, j\, i}{K}\right),\qquad
		j,i \in \{0,\ldots,K-1\}.
		\]
		Em particular, $\phi_0 \equiv 1$ é a autofunção constante associada ao
		autovalor trivial $\lambda_0 = 1$.
	\end{theorem}
	
	\begin{proof}
		Como $M$ é uma matriz circulante de permutação, satisfaz $M^K = I$,
		donde seu polinômio mínimo divide $x^K - 1$. As raízes de $x^K - 1$
		são exatamente $\lambda_j = e^{2\pi ij/K}$, $j = 0,\ldots,K-1$. As
		autofunções $\phi_j$ são os harmônicos de Fourier discretos sobre
		$\Z_K$; sua ortogonalidade decorre da relação de ortogonalidade de
		caracteres:
		\[
		\ip{\phi_j}{\phi_\ell} = \frac{1}{K}\sum_{i=0}^{K-1}
		e^{2\pi i(j-\ell)i/K} = \delta_{j\ell}.\qedhere
		\]
	\end{proof}
	
	\begin{remark}[Significado do Teorema~\ref{thm:espectro}]
		O resultado demonstra, sem qualquer hipótese sobre primos, que a
		dinâmica do motor é geometricamente equivalente a um \emph{oscilador
			harmônico discreto} de entropia nula. Esta propriedade é o ponto de
		entrada para a Conjectura de Sarnak: sistemas de entropia zero têm
		comportamento ``simples'' o suficiente para que a pseudo-aleatoriedade
		dos primos se manifeste como descorrelação espectral.
	\end{remark}
	
	% ═══════════════════════════════════════════════════════════════════════════════
	\section{Função de Colisão e Reformulação Espectral de \texorpdfstring{$\HRm$}{HR-}}
	\label{sec:colisao}
	
	\subsection{A função indicadora de colisão}
	
	\begin{definition}[Função de colisão $F$]
		\label{def:colisao}
		A observável aritmética $F \in L^2(X_N)$ que registra colisões primas
		nas janelas do motor é definida por
		\[
		F(i) := \mathbf{1}_{\PP}(e_1(i))\cdot\mathbf{1}_{\PP}(e_3(i)),
		\]
		onde $\mathbf{1}_{\PP}(x) = 1$ se $x$ é primo e $0$ caso contrário.
	\end{definition}
	
	\begin{definition}[Função de peso de von Mangoldt]
		\label{def:mangoldt}
		A versão suavizada de $F$ via a função de von Mangoldt é
		\[
		w(i) := \Lambda(e_1(i))\cdot\Lambda(e_3(i)),
		\]
		onde $\Lambda(n) = \log p$ se $n = p^k$ e $0$ caso contrário. A
		passagem de $w$ para $F$ introduz erro $O(\sqrt{K}\log K)$ proveniente
		das potências de primos com $k \geq 2$.
	\end{definition}
	
	\subsection{Decomposição espectral e reformulação de \texorpdfstring{$\HRm$}{HR-}}
	
	\begin{definition}[Transformada de Fourier de $F$]
		A projeção espectral de $F$ sobre os autovalores do motor é
		\[
		\widehat{F}(j) := \ip{F}{\phi_j} = \frac{1}{K}\sum_{i=1}^{K}
		F(i)\,e^{-2\pi i j i/K}.
		\]
	\end{definition}
	
	\begin{theorem}[Reformulação espectral de $\HRm$ — incondicional]
		\label{thm:reformulacao}
		A Hipótese Restritiva Fraca $\HRm$ para o par promovido $2M^+$ é
		equivalente a qualquer das seguintes condições:
		\begin{enumerate}[label=(\alph*)]
			\item $F(i_0) \neq 0$ para algum $i_0 \in X_N$.
			\item $\displaystyle\sum_{i=1}^{K} F(i) \geq 1$.
			\item $\widehat{F}(0) = \ip{F}{\mathbf{1}} > 0$.
		\end{enumerate}
	\end{theorem}
	
	\begin{proof}
		A equivalência (a)$\Leftrightarrow$(b) é imediata pois $F(i) \in
		\{0,1\}$. Para (b)$\Leftrightarrow$(c): como $\phi_0 \equiv 1$ e
		$\mu_N$ é uniforme,
		\[
		\widehat{F}(0) = \frac{1}{K}\sum_{i=1}^K F(i) \geq 0,
		\]
		com igualdade se e somente se $F \equiv 0$. Portanto
		$\widehat{F}(0) > 0 \Leftrightarrow \sum F(i) \geq 1$.
	\end{proof}
	
	\begin{remark}[Natureza da reformulação]
		\label{rem:tautologia}
		O Teorema~\ref{thm:reformulacao} é uma equivalência, não uma prova.
		A condição $\widehat{F}(0) > 0$ é precisamente uma reescrita de $\HRm$
		na linguagem de análise harmônica. O ganho não é lógico, mas
		\emph{linguístico}: a obstrução é agora um objeto espectral para o
		qual existem técnicas modernas — descorrelação de Sarnak, normas de
		Gowers, nilsequências — que não têm análogo direto na formulação
		original.
	\end{remark}
	
	% ═══════════════════════════════════════════════════════════════════════════════
	\section{A Obstrução Espectral e Conexões com Ferramentas Modernas}
	\label{sec:obstrucao}
	
	\subsection{Pseudo-aleatoriedade e uniformidade de Gowers}
	
	\begin{definition}[Uniformidade de Gowers $U^2$]
		Uma sequência $a: \{1,\ldots,K\} \to \C$ é \emph{$U^2$-uniforme} se
		\[
		\norm{a}_{U^2(K)}^4 := \frac{1}{K^3}\sum_{n,h_1,h_2}
		a(n)\overline{a(n+h_1)}\overline{a(n+h_2)}a(n+h_1+h_2) = o(1)
		\quad (K \to \infty).
		\]
	\end{definition}
	
	Um resultado central na combinatória aditiva moderna é que a função de
	von Mangoldt normalizada é $U^2$-uniforme ao longo de progressões
	aritméticas (Green-Tao \cite{greentao}, Gowers \cite{gowers}).
	Intuitivamente, isso significa que os primos não exibem correlações a
	dois pontos em qualquer progressão aritmética fixa — propriedade que,
	na linguagem do motor, se traduz como descorrelação entre $w(t)$ e as
	frequências $j \neq 0$ do sistema.
	
	\subsection{Descorrelação Espectral via Nilsequências}
	
	A versão anterior deste resultado invocava a Conjectura de Sarnak
	\cite{sarnak} — sobre a ortogonalidade de $\mu(n)$ a sistemas de
	entropia zero — para obter descorrelação de $\Lambda$. Essa passagem
	de $\mu$ para $\Lambda$ é não trivial e torna o resultado condicional.
	O teorema de Green-Tao-Ziegler \cite{greentaoziegler} oferece uma rota
	incondicional mais direta.
	
	\begin{proposition}[Descorrelação espectral — incondicional]
		\label{prop:sarnak}
		Seja $(X_N, \mu_N, T)$ o sistema dinâmico do motor. Para todo
		$j \neq 0$, a autofunção
		\[
		\phi_j(i) = \exp\!\left(\frac{2\pi i\,j\,i}{K}\right)
		\]
		é uma nilsequência de grau $1$ (caractere linear sobre $\Z_K$).
		Pelo teorema de Green-Tao-Ziegler \cite{greentaoziegler} sobre
		sistemas lineares em primos — especificamente, sobre a
		ortogonalidade do produto $\Lambda(L_1(i))\Lambda(L_2(i))$ a
		nilsequências para pares de formas lineares não degenerados —
		a função de peso $w(i) = \Lambda(e_1(i))\,\Lambda(e_3(i))$ satisfaz
		\[
		\lim_{K\to\infty}\widehat{w}(j) = 0,\qquad \forall\, j \neq 0.
		\]
	\end{proposition}
	
	\begin{proof}
		Consideramos o par de formas lineares
		\[
		L_1(i) = e_1(1) + 2(i-1),\qquad L_2(i) = e_3(1) - 2(i-1),
		\]
		com coeficientes $+2$ e $-2$ (Observação~\ref{rem:movimento}).
		Este par é \emph{não degenerado} — os coeficientes são linearmente
		independentes e a progressão é localmente solúvel — e satisfaz
		$L_1(i) + L_2(i) = 2M^+$ constante.
		
		Pelo teorema principal de Green-Tao-Ziegler
		\cite{greentaoziegler}, para qualquer nilsequência $\psi$ de grau finito,
		\[
		\frac{1}{K}\sum_{i=1}^{K} \Lambda(L_1(i))\,\Lambda(L_2(i))\,\psi(i) \to 0
		\quad (K \to \infty).
		\]
		Para cada $j \neq 0$, a autofunção $\phi_j(i) = e^{2\pi i j i/K}$ é
		uma nilsequência de grau $1$.
		Tomando $\psi(i) = \overline{\phi_j(i)}$, obtemos
		\[
		\widehat{w}(j)
		= \frac{1}{K}\sum_{i=1}^{K} \Lambda(e_1(i))\,\Lambda(e_3(i))\,
		\overline{\phi_j(i)} \to 0. \qedhere
		\]
	\end{proof}
	
	\begin{remark}[Relação com a Conjectura de Sarnak]
		\label{rem:sarnak}
		A Conjectura de Sarnak \cite{sarnak} — sobre $\mu(n)$ ortogonal a
		sistemas de entropia zero — é mais geral mas ainda não provada. O
		resultado acima não a utiliza: ele decorre diretamente do teorema de
		nilsequências de Green-Tao-Ziegler, que é incondicional. A Conjectura
		de Sarnak permanece como motivação conceitual para a abordagem, mas
		não como hipótese necessária.
	\end{remark}
	
	\begin{remark}[Nota sobre possível novidade na literatura]
		\label{rem:novidade}
		A aplicação do teorema de Green-Tao-Ziegler diretamente ao sistema
		de Koopman do Motor de Herança — concluindo descorrelação incondicional
		das componentes $j \neq 0$ — pode constituir uma observação nova.
		Não identificamos na literatura uma aplicação análoga de nilsequências
		a sistemas de cobertura de pares de Goldbach com soma fixa desta forma.
		Entretanto, a literatura em teoria ergódica e teoria analítica dos
		números é vasta, e uma verificação cuidadosa por especialistas é
		necessária antes de qualquer reivindicação de novidade.
	\end{remark}
	
	\subsection{A obstrução espectral e o que ela implica}
	
	A Proposição~\ref{prop:sarnak} estabelece que as componentes
	$\widehat{w}(j)$ para $j \neq 0$ tendem a zero. A componente
	$\widehat{w}(0)$, por sua vez, é a média orbital de $w$, e é
	precisamente o objeto cuja positividade equivale a $\HRm$.
	
	\begin{remark}[A barreira remanescente — formulação precisa]
		\label{rem:barreira}
		A lacuna analítica central pode ser formulada com precisão como:
		\[
		\text{provar incondicionalmente que}\quad
		\frac{1}{K}\sum_{i=1}^K \Lambda(e_1(i))\,\Lambda(e_3(i)) \gg
		\frac{1}{(\log K)^2}.
		\]
		onde $e_1(i)+e_3(i) = 2M^+$ para todo $i$, por construção
		geométrica do acoplamento $\Phi$.  Esta estimativa é equivalente
		à própria Conjectura de Goldbach na linguagem da análise harmônica.
		A reformulação não simplifica essa dificuldade; ela a expressa num
		objeto — a componente de frequência zero de $w$ no sistema de
		Koopman — para o qual as técnicas de nilsequências e médias de Weyl
		têm tração natural.
	\end{remark}
	
	\begin{remark}[Conexão com a heurística de Hardy-Littlewood]
		\label{rem:HL}
		A Conjectura de Hardy-Littlewood \cite{hardy-littlewood} prediz que
		\[
		\widehat{w}(0) \sim \mathfrak{S}(2M^-)\cdot\frac{(\log M^-)^2}{(\log K)^2} > 0,
		\]
		onde $\mathfrak{S}(2M^-)$ é a série singular, sempre positiva. Esta
		predição é consistente com a Observação~\ref{rem:barreira}: assumir
		$\widehat{w}(0) > 0$ incondicionalmente é equivalente a provar Goldbach.
		Sob GRH, os erros de Siegel-Walfisz são suficientemente controlados
		para que a cota inferior $\widehat{w}(0) > 0$ valha condicionalmente.
		Esses fatos mostram que a reformulação espectral está calibrada
		corretamente: ela captura a mesma dificuldade que a heurística clássica,
		numa linguagem que a expõe mais nitidamente.
	\end{remark}
	
	% ═══════════════════════════════════════════════════════════════════════════════
	\section{Por que os Métodos Clássicos Colapsam}
	\label{sec:crivo}
	
	Esta seção formaliza o diagnóstico mencionado na Introdução, descrevendo
	com precisão as camadas onde a abordagem clássica falha. O objetivo não é
	apenas registrar o colapso, mas mostrar por que a formulação espectral
	evita esses obstáculos específicos — ainda que não os resolva.
	
	\subsection{Configuração e notação}
	
	Fixamos $2M^-$ suficientemente grande e $L = \lfloor 2(\log M^-)^2 \rfloor$.
	Definimos
	\[
	\mathcal{N}(L, 2M^-) = \#\{p \leq L : p \in \PP,\; 2M^- - p \in \PP\}.
	\]
	A condição $\mathcal{N}(L, 2M^-) \geq 1$ é exatamente $\HRm$ com a
	progressão $\{3, 5, \ldots, L\}$, $T \sim (\log M^-)^2$.
	
	\subsection{Colapso do método do círculo}
	
	As somas de von Mangoldt $f(\alpha) = \sum_{j=1}^T \Lambda(2j+1)\,e(\alpha(2j+1))$
	e $g(\alpha)$ admitem a identidade de Fourier
	\[
	\sum_{j=1}^T \Lambda(2j+1)\,\Lambda(2M^- - (2j+1)) = \int_0^1 f(\alpha)\,\overline{g(\alpha)}\,d\alpha.
	\]
	Nos arcos maiores $\mathfrak{M}(a,q)$ com $q \leq Q = (\log M^-)^B$, o
	produto $f\cdot\overline{g}$ contém o fator $e(2M^-\beta)$ com
	$|\beta| \leq Q/(qT)$. A integral
	\[
	\int_{-Q/T}^{Q/T}|F(\beta,T)|^2\,e(2M^-\beta)\,d\beta
	= O\!\left(\frac{T^2\log T}{M^{-}Q}\right) \to 0
	\]
	quando $M^- \to \infty$ com $T = O((\log M^-)^2)$ fixo. O termo
	principal é destruído pelas oscilações rápidas de $e(2M^-\alpha)$
	à escala $M^- \gg T$.
	
	\subsection{Colapso do crivo de Rosser-Iwaniec}
	
	Crivando $\mathcal{A} = \{2M^- - p : p \leq L\}$ até nível
	$D = (\log M^-)^{1-\varepsilon}$, o erro acumulado do crivo é
	\[
	\mathcal{R} \leq \sum_{\substack{d \leq D \\ d \mid P(w)}} |r_d|
	\ll D \cdot \frac{L}{(\log L)^A}
	\sim \frac{(\log M^-)^{3-\varepsilon}}{(\log\log M^-)^A},
	\]
	enquanto o termo principal do crivo linear é apenas
	$O((\log M^-)^2/(\log\log M^-)^2)$. Como $3 - \varepsilon > 2$ para
	$\varepsilon < 1$, o resto $\mathcal{R}$ domina o sinal e a cota
	inferior colapsa.
	
	\begin{center}
		\renewcommand{\arraystretch}{1.4}
		\begin{tabular}{llll}
			\toprule
			\textbf{Camada} & \textbf{Ferramenta} & \textbf{Resultado} & \textbf{Status} \\
			\midrule
			1 & Método do círculo & Integral $\to 0$ para $M^- \gg T$ & Colapsa \\
			2 & Bombieri-Vinogradov & Inaplicável: $L \ll M^-$ & Inaplicável \\
			3 & Crivo linear (Iwaniec) & $S(\mathcal{A},z)$ positivo; sozinho não fecha & Parcial \\
			4 & Buchstab + Chen & $\mathcal{R} \gg S(\mathcal{A},z)$ & Colapsa \\
			\bottomrule
		\end{tabular}
	\end{center}
	
	\subsection{O que a abordagem espectral reposiciona}
	
	A formulação espectral não elimina a dificuldade de provar
	$\widehat{F}(0) > 0$ incondicionalmente. O que ela oferece é:
	\begin{enumerate}[label=(\alph*)]
		\item Um \emph{diagnóstico preciso} da obstrução: a dificuldade não
		está espalhada pelo problema, mas concentrada na componente de frequência
		zero de um único objeto espectral.
		\item Uma \emph{linguagem alinhada} com as técnicas que resolveram
		problemas análogos: Green-Tao usou exatamente a transição
		``contagem $\to$ fluxo'' para progressões aritméticas em primos;
		Helfgott \cite{helfgott} usou análise de Fourier refinada para
		Goldbach ternário.
		\item Uma \emph{separação clara} entre o que está provado (a estrutura
		dinâmica do motor) e o que permanece aberto (a positividade de
		$\widehat{F}(0)$).
	\end{enumerate}
	
	% ═══════════════════════════════════════════════════════════════════════════════
	\section{Síntese, Alcance e Trabalhos Futuros}
	\label{sec:conclusao}
	
	\subsection{O que este artigo prova}
	
	\begin{center}
		\renewcommand{\arraystretch}{1.4}
		\begin{tabular}{lll}
			\toprule
			\textbf{Resultado} & \textbf{Status} & \textbf{Referência} \\
			\midrule
			Sistema $(X_N,\mu_N,T)$ bem definido & Incondicional & Def.~\ref{def:estados}--\ref{def:koopman} \\
			$T$ ergódico, $h(T) = 0$ & Incondicional & Lema~\ref{lem:ergodico} \\
			Espectro discreto de $U_T$ (raízes da unidade) & Incondicional & Teo.~\ref{thm:espectro} \\
			$\HRm \Leftrightarrow \widehat{F}(0) > 0$ & Incondicional & Teo.~\ref{thm:reformulacao} \\
			Descorrelação $\widehat{w}(j) \to 0$ para $j \neq 0$ & Incondicional (GTZ) & Prop.~\ref{prop:sarnak} \\
			$\widehat{w}(0) > 0$ & Cond.\ (H-L ou GRH) & Obs.~\ref{rem:HL} \\
			\midrule
			\textbf{$\HRm$ provada} & \textbf{Não provada} & Obs.~\ref{rem:barreira} \\
			\bottomrule
		\end{tabular}
	\end{center}
	
	A última linha da tabela é o ponto central: o artigo reformula $\HRm$,
	não a demonstra. A equivalência $\HRm \Leftrightarrow \widehat{F}(0) > 0$
	é uma tautologia precisa cuja utilidade está em tornar a barreira analítica
	explícita e acessível às ferramentas modernas.
	
	\subsection{O que a reformulação oferece à comunidade}
	
	\begin{enumerate}
		\item \textbf{Localização da obstrução.} A Conjectura de Goldbach para
		$2M^+$ equivale à positividade de uma única componente de Fourier
		$\widehat{F}(0)$. Isso não simplifica o problema, mas o
		\emph{concentra}: em vez de uma afirmação difusa sobre todos os pares,
		temos uma afirmação espectral localizada.
		
		\item \textbf{Alinhamento com a escola moderna.} O Motor de Herança,
		formulado como sistema de Koopman de entropia zero, encaixa-se
		naturalmente na estrutura onde o teorema de nilsequências de
		Green-Tao-Ziegler e a Conjectura de Sarnak são aplicáveis.
		A descorrelação das componentes $j \neq 0$ já é incondicional
		(Proposição~\ref{prop:sarnak}); a única barreira remanescente
		é a componente de frequência zero.
		
		\item \textbf{Separação limpa entre camadas.} A camada algébrica
		(incondicional) e a camada aritmética (condicional) estão agora
		separadas com precisão. Qualquer progresso futuro em $\widehat{F}(0)$
		se beneficia de toda a estrutura dinâmica já estabelecida.
	\end{enumerate}
	
	\subsection{Direções futuras}
	
	\begin{enumerate}
		\item \textbf{Nilsequências e normas de Gowers de ordem superior.} A
		decomposição de $F$ em contribuições de nilsequências (no espírito de
		Green-Tao-Ziegler \cite{greentao}) pode fornecer controle sobre
		$\widehat{F}(0)$, se a progressão $\{e_1(t)\}_{t=1}^K$ exibir
		estrutura de nilsequência.
		
		\item \textbf{Prova condicional sob GRH.} Sob a Hipótese de Riemann
		Generalizada, os erros de Siegel-Walfisz satisfazem
		$|r_d| \ll d^{1/2}(\log M^-)^2$, e $\mathcal{R} \ll
		D^{3/2}(\log M^-)^2$. Para $D = (\log M^-)^{1-\varepsilon}$ com
		$\varepsilon$ próximo de 1, isso tornaria $\mathcal{R} \ll
		S(\mathcal{A},w)$ e fecharia a cota inferior $\mathcal{N} \geq 1$
		condicionalmente.
		
		\item \textbf{Recorrência de Furstenberg.} A ergodicidade de $T$ e
		a densidade positiva de $A_N$ (sob hipótese de Hardy-Littlewood)
		sugerem uma prova de recorrência múltipla: a órbita de qualquer ponto
		em $X_N$ deveria visitar $A_N$ no horizonte $K = N-1$. Formalizar
		isso via teoremas de recorrência de Furstenberg-Katznelson é uma
		direção concreta.
		
		\item \textbf{Reformulação com $L = M^\delta$.} Substituindo
		$L = (\log M^-)^2$ por $L = (M^-)^\delta$ para $\delta > 0$ fixo,
		$|\mathcal{A}|$ cresce polinomialmente e $\mathcal{R} \ll
		|\mathcal{A}|/(\log |\mathcal{A}|)^{A-1}$ por Siegel-Walfisz padrão.
		Esta versão mais forte de $\HRm$ pode ser demonstrável
		incondicionalmente, mas requer verificar que ela permanece mais fraca
		do que Goldbach.
	\end{enumerate}
	
	% ── Referências ───────────────────────────────────────────────────────────────
	\begin{thebibliography}{99}
		
		\bibitem{goldbach1}
		T.~Bandeira,
		\textit{Uma Abordagem Unificada para a Conjectura de Goldbach: Estrutura
			Combinatória, Saturação Geométrica e o Elo entre Trios e Pares Primos},
		Preprint, maio de 2026.
		
		\bibitem{grade}
		T.~Bandeira,
		\textit{Uma propriedade de soma constante na grade $3 \times N$ e as
			conjecturas de Goldbach},
		Preprint, maio de 2026.
		
		\bibitem{ancora}
		T.~Bandeira,
		\textit{O Invariante Âncora da Grade $G_{3,N}$ e um Estimador Geométrico
			Oracle-Free para a Conjectura de Goldbach},
		Preprint, maio de 2026.
		
		\bibitem{crivo_b}
		T.~Bandeira,
		\textit{Problema B do Crivo Geométrico: Demonstração da Desigualdade
			Estrutural},
		Preprint, maio de 2026.
		
		\bibitem{motor}
		T.~Bandeira,
		\textit{Motor de Herança Estrutural: Formalização Algébrica da Fita-Dobra
			e do Mecanismo de Promoção de Pares de Goldbach},
		Preprint, maio de 2026.
		
		\bibitem{greentao}
		B.~Green e T.~Tao,
		\textit{The primes contain arbitrarily long arithmetic progressions},
		Annals of Mathematics, vol.~167, pp.~481--547, 2008.
		
		\bibitem{greentaoziegler}
		B.~Green, T.~Tao e T.~Ziegler,
		\textit{An inverse theorem for the Gowers $U^{s+1}[N]$-norm},
		Annals of Mathematics, vol.~176, pp.~1231--1372, 2012.
		
		\bibitem{gowers}
		W.~T.~Gowers,
		\textit{A new proof of Szemerédi's theorem},
		Geometric and Functional Analysis, vol.~11, pp.~465--588, 2001.
		
		\bibitem{sarnak}
		P.~Sarnak,
		\textit{Three lectures on the Möbius function, randomness and dynamics},
		Publications mathématiques de l'IHÉS, 2011.
		
		\bibitem{helfgott}
		H.~A.~Helfgott,
		\textit{Major arcs for Goldbach's theorem},
		arXiv:1305.2897, 2013.
		
		\bibitem{hardy-littlewood}
		G.~H.~Hardy e J.~E.~Littlewood,
		\textit{Some problems of `Partitio Numerorum'; III: On the expression of a
			number as a sum of primes},
		Acta Mathematica, vol.~44, pp.~1--70, 1923.
		
		\bibitem{iwaniec}
		H.~Iwaniec e E.~Kowalski,
		\textit{Analytic Number Theory},
		American Mathematical Society Colloquium Publications, vol.~53, 2004.
		
		\bibitem{chen}
		J.~R.~Chen,
		\textit{On the representation of a large even integer as the sum of a prime
			and the product of at most two primes},
		Scientia Sinica, vol.~16, pp.~157--176, 1973.
		
	\end{thebibliography}
	
\end{document}